%% front/how-to-read.tex -- explains the fixed entry format to the reader.
\chapter*{How to Read an Entry}
\hbtopicmark{HOW TO READ AN ENTRY}

The handbook is organised by \emph{subject}, not by book. Sources are merged
into a single argument, so a reader never has to know which book said what in
order to follow the text. Provenance is available on demand, not in the way.

\section*{The anatomy of an entry}

\begingroup\small
\begin{itemize}
\item \textbf{Core} -- the whole entry in one to three sentences. If you read
  nothing else, read this. It is boxed because it must never be buried.
\item \textbf{Mechanism} -- why the move works on a human nervous system.
  A tactic whose mechanism you understand can be adapted; one you merely
  memorise breaks on first contact.
\item \textbf{Tells} -- how to spot it being used, usually on you. Bulleted,
  observable, and deliberately behavioural rather than intuitive.
\item \textbf{Moves} -- how it is executed. Numbered and imperative. This is
  the slot that makes the entry useful defensively: you cannot recognise a
  sequence you have never seen described.
\item \textbf{Counters} -- what to do about it. Given equal weight to Moves.
\item \textbf{Cost} -- when the move fails, what it costs the user, and the
  blowback it produces. Every technique in this book has a failure mode, and
  an entry without one is incomplete rather than optimistic.
\item \textbf{Field note} -- a short modern illustration written for this
  handbook. Not from the sources.
\end{itemize}
\par\endgroup

\section*{Notation}

\begingroup\small
\begin{itemize}
\item \texttt{[B1]}, \texttt{[B3]} -- superscript tags naming the source a
  claim came from. Full details in Appendix~A and \texttt{registry/books.yaml}.
\item \textbf{Drawn from:} -- the entry's complete source list.
\item \textbf{Related:} -- neighbouring entries, named in full. The book is
  a lattice, not a line; follow the links when an entry assumes one. A
  link marked \emph{planned} points at an entry reserved in the registry
  but not yet written.
\item \texttt{T-\,nn-\,nn} -- entry number: chapter, then position in
  chapter. It is printed beside every heading so citations and the source
  registry stay precise; you never need it to read.
\end{itemize}
\par\endgroup

\section*{Order of reading}

Parts~I and~II are prerequisites: they cover what manipulation is, who does
it, and how to read a person. Parts~III--V are the technique layers, and can
be read in any order. Part~VI is defence and can be read first by anyone who
is currently being worked on. Part~VII is the cost accounting.
